% Libro 2 · Mundo 5 · Probabilidad (app: pu-probabilidad)
\mundo{Probabilidad}{pu-probabilidad}{La probabilidad mide qué tan esperable es algo, con un número entre 0 (imposible) y
1 (seguro). En el examen es casi siempre conteo disfrazado: casos favorables sobre casos posibles. La dificultad está en
contar bien los posibles y en no confundir «y» con «o».}{%
\item \textbf{5.1} Espacio muestral y regla de Laplace
\item \textbf{5.2} Complemento, independencia y sin reposición}

\leccion{pu-pb-1}{Espacio muestral y Laplace}{Juego de la tabla de los dos dados.}

\begin{reto}
Se lanzan dos dados. ¿Cuál es la probabilidad de que la suma sea 8?
\opciones{$\tfrac{5}{36}$}{$\tfrac16$}{$\tfrac{8}{36}$}{$\tfrac19$}
\end{reto}

\subsection{La idea: todos los casos con el mismo peso}
El \clave{espacio muestral} es la lista de todos los resultados posibles de un experimento. Si todos son
\emph{igualmente probables}, la probabilidad de un evento es la fracción de resultados que lo cumplen:
\[P(E) = \frac{\text{casos favorables}}{\text{casos posibles}}\qquad\text{(regla de Laplace).}\]

\subsubsection{Dos dados: 36 casos, no 11}
Las sumas posibles de dos dados van de 2 a 12 (11 valores), pero \emph{no} son igualmente probables. Los casos que sí lo
son son los 36 pares (dado 1, dado 2):

\begin{center}\small
\begin{tabular}{c|cccccc}
 & 1 & 2 & 3 & 4 & 5 & 6\\\hline
1 & 2 & 3 & 4 & 5 & 6 & 7\\
2 & 3 & 4 & 5 & 6 & 7 & \cellcolor{zlime!50}8\\
3 & 4 & 5 & 6 & 7 & \cellcolor{zlime!50}8 & 9\\
4 & 5 & 6 & 7 & \cellcolor{zlime!50}8 & 9 & 10\\
5 & 6 & 7 & \cellcolor{zlime!50}8 & 9 & 10 & 11\\
6 & 7 & \cellcolor{zlime!50}8 & 9 & 10 & 11 & 12\\
\end{tabular}
\end{center}

La suma 7 es la más frecuente (6 casos); 2 y 12 las menos (1 caso cada una).

\begin{regla}
\begin{itemize}
  \item $0 \le P(E) \le 1$. Imposible: 0. Seguro: 1.
  \item Laplace solo vale si los casos son igualmente probables: cuenta los casos «finos» (pares ordenados, cartas
  individuales), no los resultados agrupados.
  \item Si dos eventos no pueden ocurrir a la vez, $P(A \text{ o } B) = P(A) + P(B)$.
\end{itemize}
\end{regla}

\subsection{Ejemplos resueltos}
\begin{ejemplo}[Una bolsa]
Una bolsa tiene 4 bolas rojas, 5 azules y 3 verdes. Se saca una. ¿Probabilidad de que no sea azul?
\solucion No azules: 7 de 12. $P = \tfrac{7}{12}$.
\end{ejemplo}

\begin{ejemplo}[Tres monedas]
Se lanzan tres monedas. ¿Probabilidad de obtener exactamente dos caras?
\solucion Espacio: $2^3 = 8$ secuencias. Con dos caras: CCS, CSC, SCC. $P = \tfrac38$.
\end{ejemplo}

\begin{ejemplo}[Un número al azar]
Se elige al azar un número del 1 al 50. ¿Probabilidad de que sea múltiplo de 3 o de 5?
\solucion Múltiplos de 3: 16; de 5: 10; de 15: 3. Favorables: $16 + 10 - 3 = 23$. $P = \tfrac{23}{50}$.
\end{ejemplo}

\begin{trampa}
Usar las 11 sumas posibles de dos dados como si fueran igualmente probables: «$P(\text{suma } 8) = \tfrac{1}{11}$» es
falso.
\end{trampa}

\taller
\piso{Entrena}{contar casos}
\begin{ejercicios}
\item Con dos dados, $P(\text{suma} = 7)$. \sol{$\tfrac{6}{36} = \tfrac16$.}
\item Con dos dados, $P(\text{dobles})$. \sol{$\tfrac16$.}
\item Se elige una letra de la palabra MATEMÁTICA. $P(\text{vocal})$. \sol{A, E, A, I, A: 5 de 10: $\tfrac12$.}
\end{ejercicios}
\piso{Como en el examen}{opción múltiple}
\begin{ejercicios}
\item Se lanzan dos dados. La probabilidad de que el producto sea par es
\opciones{$\tfrac12$}{$\tfrac34$}{$\tfrac14$}{$\tfrac23$}
\sol{B. Impar solo si los dos son impares: $\tfrac{9}{36} = \tfrac14$; par: $\tfrac34$.}
\item Una caja tiene 20 fichas numeradas del 1 al 20. ¿Probabilidad de sacar un número primo?
\opciones{$\tfrac{2}{5}$}{$\tfrac{9}{20}$}{$\tfrac{1}{2}$}{$\tfrac{7}{20}$}
\sol{A. 2, 3, 5, 7, 11, 13, 17, 19: 8 de 20.}
\end{ejercicios}

\begin{retoresuelto}
Suma 8: (2,6), (3,5), (4,4), (5,3), (6,2): 5 casos de 36. $P = \tfrac{5}{36}$. \textbf{Opción A.}
\end{retoresuelto}
\bitacora

\leccion{pu-pb-2}{Complemento e independencia}{Juego del árbol de probabilidades.}

\begin{reto}
Una urna tiene 3 bolas rojas y 2 azules. Se sacan dos, una tras otra, sin devolver la primera. ¿Probabilidad de que las
dos sean rojas?
\opciones{$\tfrac{9}{25}$}{$\tfrac{3}{10}$}{$\tfrac35$}{$\tfrac12$}
\end{reto}

\subsection{La idea: el árbol multiplica}
Cuando un experimento ocurre en etapas, la probabilidad de un camino del árbol es el \clave{producto} de las
probabilidades de sus ramas. Si las etapas no se afectan (lanzar una moneda y luego un dado), los eventos son
\clave{independientes} y se multiplica directamente. Si una etapa cambia la siguiente (sacar sin devolver), la segunda
rama se calcula con lo que quedó.

\figura{arbolprob}{Sin reposición: la segunda rama se calcula con lo que quedó en la urna.}

\subsubsection{Complemento}
$P(\text{no } E) = 1 - P(E)$. Es la herramienta para «al menos uno»: la probabilidad de al menos un 6 en tres lanzamientos
es $1 - \left(\tfrac56\right)^3 = 1 - \tfrac{125}{216} = \tfrac{91}{216}$.

\begin{regla}
\begin{itemize}
  \item Independientes: $P(A \text{ y } B) = P(A)\cdot P(B)$.
  \item Sin reposición: la segunda probabilidad se calcula con lo que quedó.
  \item «Al menos uno»: $1 - P(\text{ninguno})$.
  \item Varios caminos del árbol que llevan al mismo evento: se suman.
\end{itemize}
\end{regla}

\subsection{Ejemplos resueltos}
\begin{ejemplo}[Un color de cada uno]
Con la urna del reto, ¿probabilidad de sacar una de cada color?
\solucion Caminos RA y AR: $\tfrac35 \cdot \tfrac24 + \tfrac25 \cdot \tfrac34 = \tfrac{6}{20} + \tfrac{6}{20} = \tfrac35$.
\end{ejemplo}

\begin{ejemplo}[Examen al azar]
Un estudiante responde al azar 3 preguntas de 4 opciones. ¿Probabilidad de acertar al menos una?
\solucion Fallar una: $\tfrac34$. Fallar las tres: $\tfrac{27}{64}$. Al menos un acierto: $\tfrac{37}{64}$.
\end{ejemplo}

\begin{trampa}
Sumar cuando hay que multiplicar: «$P(\text{dos rojas}) = \tfrac35 + \tfrac24$» da más de 1, lo cual es imposible. Una
probabilidad nunca supera 1.
\end{trampa}

\taller
\piso{Entrena}{árboles}
\begin{ejercicios}
\item Se lanza una moneda y un dado. $P(\text{cara y número mayor que 4})$. \sol{$\tfrac12 \cdot \tfrac26 = \tfrac16$.}
\item De una baraja de 40 cartas (10 de cada palo) se sacan dos sin reposición. $P(\text{ambas de oros})$. \sol{$\tfrac{10}{40} \cdot \tfrac{9}{39} = \tfrac{3}{52}$.}
\item Se lanza una moneda 4 veces. $P(\text{al menos una cara})$. \sol{$1 - \tfrac{1}{16} = \tfrac{15}{16}$.}
\end{ejercicios}
\piso{Como en el examen}{opción múltiple}
\begin{ejercicios}
\item La probabilidad de que llueva el sábado es $0.3$ y el domingo $0.4$, de forma independiente. ¿Probabilidad de que no llueva ningún día?
\opciones{$0.42$}{$0.3$}{$0.58$}{$0.12$}
\sol{A. $0.7 \cdot 0.6$.}
\item Una caja tiene 2 bombillos dañados y 3 buenos. Se prueban uno a uno sin devolverlos. ¿Probabilidad de que los dos primeros sean dañados?
\opciones{$\tfrac{1}{10}$}{$\tfrac{4}{25}$}{$\tfrac15$}{$\tfrac25$}
\sol{A. $\tfrac25 \cdot \tfrac14$.}
\end{ejercicios}

\begin{retoresuelto}
Primera roja: $\tfrac35$. Queda una urna con 2 rojas y 2 azules: segunda roja $\tfrac24$. Producto: $\tfrac{6}{20} =
\tfrac{3}{10}$. \textbf{Opción B.} La A sería con reposición.
\end{retoresuelto}
\bitacora

\evaluacion
\begin{ejercicios}
\item Con dos dados, la probabilidad de que la suma sea mayor que 9 es
\opciones{$\tfrac16$}{$\tfrac{1}{12}$}{$\tfrac{1}{4}$}{$\tfrac{5}{36}$}
\sol{A. Sumas 10, 11, 12: $3 + 2 + 1 = 6$ casos. (5.1)}
\item De un grupo de 12 hombres y 8 mujeres se escoge una persona al azar. $P(\text{mujer})$ es
\opciones{$\tfrac25$}{$\tfrac23$}{$\tfrac{8}{12}$}{$\tfrac35$}
\sol{A. (5.1)}
\item Se lanzan dos monedas. La probabilidad de que salgan dos caras es
\opciones{$\tfrac13$}{$\tfrac14$}{$\tfrac12$}{$\tfrac34$}
\sol{B. (5.1)}
\item Se lanza un dado tres veces. La probabilidad de que nunca salga 1 es
\opciones{$\tfrac{125}{216}$}{$\tfrac{91}{216}$}{$\tfrac12$}{$\tfrac{5}{18}$}
\sol{A. (5.2)}
\item Una urna tiene 4 blancas y 1 negra. Se sacan dos sin reposición. $P(\text{la negra sale})$ es
\opciones{$\tfrac25$}{$\tfrac15$}{$\tfrac{9}{25}$}{$\tfrac12$}
\sol{A. $1 - \tfrac45 \cdot \tfrac34 = 1 - \tfrac35$. (5.2)}
\end{ejercicios}
